\documentclass[a4paper]{article}
% Español (México) con XeLaTeX: fontspec para fuentes Unicode, polyglossia
% para la división silábica y los nombres del documento en la variante mexicana.
\usepackage{fontspec}
\usepackage{polyglossia}
\setdefaultlanguage[variant=mexican]{spanish}
% Los nombres chinos van como «pinyin (original)»: xeCJK da a los caracteres
% Han su propia fuente; sin ella XeLaTeX los omite con un aviso.
% FandolSong no cubre algunos caracteres de nombres (U+73FA); WenQuanYi Zen Hei sí.
\usepackage[AutoFallBack=true]{xeCJK}
\setCJKmainfont{FandolSong-Regular.otf}
\setCJKfallbackfamilyfont{\CJKrmdefault}{WenQuanYi Zen Hei}
% polyglossia no traduce los nombres de listings.
\AtBeginDocument{\providecommand{\lstlistingname}{}\renewcommand{\lstlistingname}{Listado}%
  \providecommand{\lstlistlistingname}{}\renewcommand{\lstlistlistingname}{Índice de listados}}
\usepackage{amsmath, amssymb}
\usepackage{graphicx}
\usepackage[margin=2.5cm]{geometry}
\usepackage[most]{tcolorbox}
\usepackage{etoolbox}
\usepackage{listings}
\usepackage{hyperref}
\usepackage{booktabs}
\usepackage{float}
\newtcolorbox{knowledgebox}[1]{enhanced, colback=blue!5!white, colframe=blue!75!black, colbacktitle=blue!75!black, coltitle=white, fonttitle=\bfseries, title={#1}, attach boxed title to top left={yshift=-2mm, xshift=2mm}, boxrule=1pt, sharp corners}
\newtcolorbox{importantbox}[1]{enhanced, colback=yellow!10!white, colframe=yellow!80!black, colbacktitle=yellow!80!black, coltitle=black, fonttitle=\bfseries, title={#1}, sharp corners}
\newtcolorbox{warningbox}[1]{enhanced, colback=red!5!white, colframe=red!75!black, colbacktitle=red!75!black, coltitle=white, fonttitle=\bfseries, title={#1}, sharp corners}
\lstset{language=Python, basicstyle=\ttfamily\small, keywordstyle=\color{blue}, stringstyle=\color{red!60!black}, commentstyle=\color{green!60!black}, breaklines=true, frame=single, numbers=left, numberstyle=\tiny\color{gray}, captionpos=b, extendedchars=true}
\newcommand{\notetitle}{veRL Agentic Loop: explicación detallada del código}
\newcommand{\noteauthors}{Elaboradas a partir de materiales públicos del curso}
\newcommand{\notedate}{2025}
\newcommand{\videochannel}{Wudaokou Nash (五道口纳什)}
\newcommand{\videopublishdate}{2025}
\newcommand{\videoduration}{aproximadamente 20 minutos}
\newcommand{\videourl}{}
\newcommand{\videocoverpath}{cover.jpg}
\newcommand{\repourl}{https://github.com/hqhq1025/ai-course-notes}

% Footer with repo info
\usepackage{fancyhdr}
\pagestyle{fancy}
\fancyhf{}
\fancyfoot[C]{\thepage}
\fancyfoot[R]{\footnotesize\color{gray}\href{\repourl}{AI Course Notes}}
\renewcommand{\headrulewidth}{0pt}
\renewcommand{\footrulewidth}{0pt}

\begin{document}
\begin{titlepage}\centering
{\Large Notas del curso\par}\vspace{1.2cm}{\huge\bfseries \notetitle\par}\vspace{0.8cm}{\large \noteauthors\par}\vspace{0.3cm}{\large \notedate\par}\vspace{1.2cm}
\ifdefempty{\videocoverpath}{}{\includegraphics[width=0.82\textwidth,height=0.45\textheight,keepaspectratio]{\videocoverpath}\par}
\vfill
\begin{tcolorbox}[width=0.9\textwidth, colback=black!2!white, colframe=black!60, sharp corners]
\textbf{Autor o canal del video}: \videochannel\par\textbf{Fecha de publicación}: \videopublishdate\par\textbf{Duración del video}: \videoduration\par
\ifdefempty{\videourl}{}{\textbf{Enlace del video}: \href{\videourl}{\nolinkurl{\videourl}}\par}\par
\textbf{Página del proyecto}: \href{\repourl}{\nolinkurl{\repourl}}\par
\end{tcolorbox}
\end{titlepage}
\tableofcontents\newpage

\section{Introducción}
En esta ocasión profundizamos en el código del Agentic Loop de veRL, con un análisis detallado de la implementación de \texttt{agentloop.py}.

\section{Estructura de las clases principales}
\subsection{AgentLoopBase}
Clase base abstracta; el método \texttt{run} es un método abstracto que el usuario debe personalizar. La versión oficial ofrece dos implementaciones: Singleton y TwoAgents.

\subsection{AgentLoopWorker}
Encapsula la lógica de ejecución de un solo Agent y gestiona las transiciones de estado.

\subsection{AgentLoopManager}
La clase más central; mediante el método \texttt{generate\_sequence} coordina todos los componentes:
\begin{itemize}
  \item Gestiona el server de inferencia
  \item Coordina el estado de los Worker
  \item Maneja la ejecución asíncrona de los tool call
\end{itemize}

\section{Filosofía de diseño}
El Agentic Loop de veRL busca:
\begin{itemize}
  \item Abstracciones simples
  \item Ejecución asíncrona eficiente
  \item Interfaces flexibles y personalizables
\end{itemize}

\subsection{Resumen de la sección}
AgentLoopManager es el cerebro del Agentic RL de veRL: coordina de manera unificada la inferencia, las llamadas a herramientas y el entrenamiento.

\newpage
\section{Implicaciones de diseño: el nivel de abstracción determina la escalabilidad}
Separar las tres capas de Base, Worker y Manager no es para que «el código se vea mejor», sino para que el Agentic Loop se pueda reutilizar entre distintas tareas. Lo que cambia con más rapidez es la lógica concreta de la tarea, no la gestión de estado ni la orquestación de la inferencia; por eso una capa de abstracción unificada reduce de forma notable el costo de los experimentos posteriores.

\begin{knowledgebox}{Qué revisar primero al leer el código de un framework}
\begin{itemize}
  \item Qué capa se encarga de la orchestration unificada
  \item Qué capa asume el estado concreto de la tarea
  \item Qué interfaces existen para que quienes investigan sustituyan la lógica de negocio
\end{itemize}
\end{knowledgebox}

\subsection{Resumen de la sección}
La mejor manera de entender AgentLoopManager no es memorizar nombres de funciones, sino distinguir con claridad qué responsabilidades se abstrajeron de forma estable y cuáles se dejaron deliberadamente para la lógica personalizada.

\newpage
\section{Síntesis y ampliación}
\begin{enumerate}
  \item Jerarquía AgentLoopBase $\to$ Worker $\to$ Manager
  \item generate\_sequence es el método de coordinación central
  \item Un Agent personalizado solo necesita implementar AgentLoopBase.run
\end{enumerate}
\end{document}
